\subsection{Modules of fractions of tensor products and homomorphism modules}

PROPOSITION 18. Let A be a ring and S a multiplicative subset of A.

\begin{enumerate}
\def\labelenumi{(\roman{enumi})}
\tightlist
\item
  If M and N are two A-modules, the \(S^{-1}A\) -modules ( \(S^{-1}M\) ) \(\otimes_{A}N\) , \(M\otimes_{A}(S^{-1}N)\) , ( \(S^{-1}M\) ) \(\otimes_{S^{-1}A}(S^{-1}N)\) and \(S^{-1}(M\otimes_{A}N)\) are canonically isomorphic.
\item
  If \(M'\) and \(N'\) are two \(S^{-1}A\) -modules, the canonical homomorphism

\[
\mathbf {M} ^ {\prime} \otimes_ {\mathbf {A}} \mathbf {N} ^ {\prime} \rightarrow \mathbf {M} ^ {\prime} \otimes_ {\mathbf {S} ^ {- 1} \mathbf {A}} \mathbf {N} ^ {\prime}
\]

derived from the A-bilinear mapping \((x', y') \mapsto x' \otimes y'\) from \(M' \times N'\) to \(M' \otimes_{S^{-1}} A N'\) is bijective.

\end{enumerate}
Assertion (i) is an immediate consequence of the definition \(\mathbf{S}^{-1}\mathbf{M} =\)

\(M \otimes_{A} S^{-1} A\) and the associativity of tensor products, which gives to within canonical isomorphisms

\[
\begin{array}{r l} \left(\mathrm{S} ^ {- 1} \mathrm{M} \otimes_ {\mathrm{S} ^ {- 1} \mathrm{A}} \left(\mathrm{S} ^ {- 1} \mathrm{N}\right) = (\mathrm{S} ^ {- 1} \mathrm{M}) \otimes_ {\mathrm{S} ^ {- 1} \mathrm{A}} \left(\mathrm{S} ^ {- 1} \mathrm{A} \otimes \mathrm{N}\right) = (\mathrm{S} ^ {- 1} \mathrm{M}) \otimes_ {\mathrm{A}} \mathrm{N} \right. \\ = (\mathrm{S} ^ {- 1} \mathrm{A}) \otimes_ {\mathrm{A}} \mathrm{M} \otimes_ {\mathrm{A}} \mathrm{N} = \mathrm{S} ^ {- 1} (\mathrm{M} \otimes_ {\mathrm{A}} \mathrm{N}). \end{array}
\]

To prove (ii), we note first that in \(\mathbf{M}'\) and \(\mathbf{N}'\) , considered as A-modules, the homotheties induced by the elements \(s \in S\) are bijective, hence \(\mathbf{M}' = \mathbf{S}^{-1}\mathbf{M}'\) and \(\mathbf{N}' = \mathbf{S}^{-1}\mathbf{N}'\) (no. 2, Remark 4) and similarly \(\mathbf{S}^{-1}(\mathbf{M}' \otimes_{\mathbf{A}} \mathbf{N}') = \mathbf{M}' \otimes_{\mathbf{A}} \mathbf{N}'\) ; (ii) is then a special case of (i).

COROLLARY. Let M be an A-module and a an ideal of A. The sub- \(\mathbf{S}^{-1}\mathbf{A}\) -modules \((\mathbf{S}^{-1}\mathfrak{a})(\mathbf{S}^{-1}\mathbf{M})\) , \(\mathfrak{a}(\mathbf{S}^{-1}\mathbf{M})\) , \((\mathbf{S}^{-1}\mathfrak{a})j(\mathbf{M})\) (where \(j: \mathbf{M} \to \mathbf{S}^{-1}\mathbf{M}\) is the canonical mapping) and \(\mathbf{S}^{-1}(\mathfrak{a}\mathbf{M})\) of \(\mathbf{S}^{-1}\mathbf{M}\) are identical. In particular, if \(\mathfrak{a}\) and \(\mathfrak{b}\) are two ideals of A, then

\[
(S ^ {- 1} a) (S ^ {- 1} b) = a (S ^ {- 1} b) = (S ^ {- 1} a) b = S ^ {- 1} (a b).
\]

Remark. Let M, N, P be three A-modules, \(f: M \times N \to P\) an A-bilinear mapping and \(S^{-1}f: (S^{-1}M) \times (S^{-1}N) \to (S^{-1}P)\) the \(S^{-1}A\) -bilinear mapping obtained from f by extending the ring of scalars to \(S^{-1}A\) (Algebra, Chapter IX, § 1, no. 4, Proposition 1). Let \(i: M \to S^{-1}M\) , \(j: N \to S^{-1}N\) be the canonical homomorphisms; it follows immediately that, if Q is the sub-A-module of P generated by \(f(M \times N)\) , \(S^{-1}Q\) is the sub- \(S^{-1}A\) -module of \(S^{-1}P\) generated by \((S^{-1}f)(i(M) \times j(N))\) .

PROPOSITION 19. Let A be a ring and S a multiplicative subset of A.

\begin{enumerate}
\def\labelenumi{(\roman{enumi})}
\tightlist
\item
  If M and N are two A-modules and M is finitely generated (resp. finitely presented), the canonical homomorphism (Chapter I, § 2, no. 10, formula (10))

\[
\mathrm{S} ^ {- 1} \operatorname{Hom} _ {\mathrm{A}} (\mathrm{M}, \mathrm{N}) \rightarrow \operatorname{Hom} _ {\mathrm{S} ^ {- 1} \mathrm{A}} (\mathrm{S} ^ {- 1} \mathrm{M}, \mathrm{S} ^ {- 1} \mathrm{N})
\]

is injective (resp. bijective).

\item
  If \(M'\) , \(N'\) are two \(S^{-1}A\) -modules, the canonical bijection

\[
\operatorname{Hom} _ {\mathbf {S} ^ {- 1} \mathbf {A}} (\mathbf {M} ^ {\prime}, \mathbf {N} ^ {\prime}) \rightarrow \operatorname{Hom} _ {\mathbf {A}} (\mathbf {M} ^ {\prime}, \mathbf {N} ^ {\prime})
\]

is bijective.

\end{enumerate}
As \(S^{-1}A\) is a flat A-module, (i) is a particular case of Chapter I, §2, no. 10, Proposition 11. On the other hand, we have already remarked that every A-homomorphism of \(S^{-1}A\) -modules is necessarily an \(S^{-1}A\) -homomorphism, in the course of proving Proposition 3 of no. 2; whence (ii).

PROPOSITION 20. Let A, A' be two rings, \(\rho: \mathbf{A} \to \mathbf{A}'\) a homomorphism, S a multiplicative subset of A, \(S' = \rho(S)\) and \(\rho': S^{-1}A \to S'^{-1}A'\) the homomorphism corresponding to \(\rho\) (no. 1, Proposition 2).

\begin{enumerate}
\def\labelenumi{(\roman{enumi})}
\tightlist
\item
  For every \(\mathbf{A}'\) -module \(\mathbf{M}'\) there exists a unique \(\mathbf{S}^{-1}\mathbf{A}\) -isomorphism

\[
j \colon \mathrm{S} ^ {- 1} \rho_ {*} (\mathrm{M} ^ {\prime}) \rightarrow \rho_ {*} ^ {\prime} (\mathrm{S} ^ {- 1} \mathrm{M} ^ {\prime})
\]

such that \(j(m' / s) = m' / \rho(s)\) for all \(m' \in \mathbf{M}'\) , \(s \in \mathbf{S}\) .

\item
  For every A-module M, there exists a unique isomorphism

\[
j ^ {\prime} \colon (\mathrm{S} ^ {- 1} \mathrm{M}) \otimes_ {\mathrm{S} ^ {- 1} \mathrm{A}} (\mathrm{S} ^ {\prime - 1} \mathrm{A} ^ {\prime}) \rightarrow \mathrm{S} ^ {\prime - 1} (\mathrm{M} \otimes_ {\mathrm{A}} \mathrm{A} ^ {\prime})
\]

of \(\mathbf{S}^{\prime -1}\mathbf{A}^{\prime}\) -modules such that \(j^{\prime}((m / s)\otimes (a^{\prime} / s^{\prime})) = (m\otimes a^{\prime}) / (\rho (s)s^{\prime})\)

\end{enumerate}
\begin{enumerate}
\def\labelenumi{(\roman{enumi})}
\item
  If we consider \(S^{\prime-1}M^{\prime}\) as an A-module by means of the composite homomorphism \(i_{M^{\prime}}^{S^{\prime}}\circ\rho\) , the homotheties induced by the elements of S are bijective, hence there exists a unique homomorphism j with the stated property (no. 2, Proposition 3). As \(\rho(S)=S^{\prime}\) , j is surjective; moreover, if \(m^{\prime}\in M^{\prime}\) , \(s\in S\) , \(m^{\prime}/\rho(s)=0\) , there exists \(t^{\prime}\in S^{\prime}\) such that \(t^{\prime}m^{\prime}=0\) ; as there exists \(t\in S\) such that \(\rho(t)=t^{\prime}\) , \(t^{\prime}.m^{\prime}=0\) in \(\rho_{*}(M^{\prime})\) , hence \(m^{\prime}/s=0\) in \(S^{-1}\rho_{*}(M^{\prime})\) .
\item
  As \((\mathbf{S}^{-1}\mathbf{M})\otimes_{\mathbf{S}^{-1}\mathbf{A}}(\mathbf{S}^{\prime -1}\mathbf{A}^{\prime}) = (\mathbf{M}\otimes_{\mathbf{A}}\mathbf{S}^{-1}\mathbf{A})\otimes_{\mathbf{S}^{-1}\mathbf{A}}(\mathbf{S}^{\prime -1}\mathbf{A}^{\prime})\) and

\[
\mathbf {S} ^ {- 1} (\mathbf {M} \otimes_ {\mathbf {A}} \mathbf {A} ^ {\prime}) = (\mathbf {M} \otimes_ {\mathbf {A}} \mathbf {A} ^ {\prime}) \otimes_ {\mathbf {A} ^ {\prime}} (\mathbf {S} ^ {- 1} \mathbf {A} ^ {\prime}),
\]

the existence of \(j'\) follows from the associativity of tensor products.

\end{enumerate}
